% Capítulo 5: Pruebas y Resultados

\section{Descripción de experimentos}

Para la validación de Qsimov se han diseñado cinco experimentos, ordenados de menor a mayor
complejidad del dataset y del escenario evaluado.

\subsection{Configuración experimental común}

Todos los experimentos se ejecutan bajo las siguientes condiciones:

Todos los experimentos se ejecutaron en la GPU del servidor
\texttt{fermidi1-informat-uv-es} de la Universitat de València (NVIDIA GPU, Python 3.10). El
servidor dispone de acceso exclusivo a la GPU durante los experimentos, por lo que los tiempos de
entrenamiento no están distorsionados por la secuencialidad impuesta por el bajo número de cores de
CPU.

Se utilizó PyTorch 2.x como framework en todos los experimentos reportados. La librería Qsimov
ofrece implementaciones equivalentes en TensorFlow/Keras, pero todos los experimentos presentados en
este documento se han reejecutado con el backend PyTorch sobre GPU para garantizar comparativas de
tiempo homogéneas entre métodos. El entorno software concreto fue:
\begin{itemize}
  \item Python 3.10
  \item PyTorch 2.x con CUDA
  \item Cython 3.0 (extensiones compiladas con \texttt{python setup.py build\_ext --inplace})
  \item MLflow para tracking de experimentos
\end{itemize}

Las semillas aleatorias están fijadas en cada experimento. El tracking de experimentos se realiza
con MLflow~\cite{ref16}. El Anexo~C detalla los comandos exactos para reproducir cada experimento.

El modelo base se pre-entrena con el método estándar (Adam, cross-entropy) antes de aplicar Qsimov.

\subsection{Experimento 1: Velocidad y precisión (MNIST)}

\paragraph{Configuración} Dataset MNIST (60.000 muestras de entrenamiento, 10.000 de test; imágenes
de dígitos manuscritos $28 \times 28$ píxeles~\cite{ref17}). Red base: CNN con dos bloques
Conv2D+MaxPooling seguidos de tres capas Dense (18.474 parámetros totales). Punto de corte:
\texttt{initial\_layer = -2} (la última capa \texttt{Linear(16} $\rightarrow$ \texttt{10)} en PyTorch, ya que
ReLU y Dropout cuentan como capas separadas). Número de caminos generados: \textbf{17} (16 entradas
de \texttt{Linear(16} $\rightarrow$ \texttt{10)} más el nodo de bias).

\begin{table}[htbp]
  \centering
  \caption{Métodos comparados en el experimento de velocidad y precisión (MNIST).}
  \label{tab:metodos_mnist}
  \begin{tabular}{lp{9cm}}
    \toprule
    \textbf{Método} & \textbf{Descripción} \\
    \midrule
    \texttt{standard} (half) & Re-entrenamiento estándar PyTorch con la mitad del dataset (30.000 muestras), 10 epochs \\
    \texttt{standard} (full) & Re-entrenamiento estándar PyTorch con el dataset completo (60.000 muestras), 10 epochs \\
    \texttt{qsimov\_gradient} & QsimovGradient sobre la mitad del dataset (PathSelector), ajustado en el dataset completo, 10 epochs \\
    \texttt{qsimov\_linear}   & QsimovLinearSystem sobre la mitad del dataset (PathSelector), ajustado en el dataset completo, una sola pasada \\
    \bottomrule
  \end{tabular}
\end{table}

Se comparan dos funciones de pérdida: entropía cruzada categórica y MSE. El \texttt{qsimov\_linear}
solo se evalúa con MSE, puesto que el sistema lineal requiere que la función objetivo sea cuadrática.

\begin{table}[htbp]
  \centering
  \caption{Resultados del experimento de velocidad y precisión sobre MNIST (GPU, PyTorch, MSE).}
  \label{tab:tabla2}
  \begin{tabular}{lcc}
    \toprule
    \textbf{Método} & \textbf{Precisión test (\%)} & \textbf{Tiempo (s)} \\
    \midrule
    \texttt{standard} (half)              & 97.42 & 72.4 \\
    \texttt{standard} (full)              & 97.32 & 144.9 \\
    \texttt{qsimov\_gradient} (half$\rightarrow$full) & 97.40 & 90.5 \\
    \texttt{qsimov\_linear} (half$\rightarrow$full)   & 97.40 & 6.6 \\
    \bottomrule
  \end{tabular}
\end{table}

El sistema lineal de Qsimov alcanza una precisión equivalente al re-entrenamiento estándar
entrenado con el dataset completo, pero en un tiempo \textbf{22 veces menor} (6.6\,s frente a
144.9\,s). Esta ventaja temporal se debe a que \texttt{QsimovLinearSystem} no realiza descenso de
gradiente iterativo: en su lugar, resuelve el sistema de ecuaciones $Ax = b$ en una única pasada
sobre los datos mediante factorización QR y back-substitution. La diferencia de precisión entre
todos los métodos es inferior a 0.1 puntos porcentuales, lo que indica que la reducción del
re-entrenamiento a los caminos activos de la última capa no implica pérdida significativa de
capacidad representacional.

\subsection{Experimento 2: Robustez al learning rate (MNIST)}

Los resultados anteriores confirmaron que el sistema lineal iguala la precisión del
re-entrenamiento estándar. El siguiente experimento se centra en el módulo de gradiente y en una
debilidad práctica frecuente del re-entrenamiento estándar: la sensibilidad al learning rate, cuya
elección incorrecta puede degradar drásticamente el rendimiento.

\paragraph{Configuración} Mismo dataset y red que el experimento anterior (PyTorch, GPU). Se evalúa
el rendimiento de \texttt{qsimov\_gradient} y del re-entrenamiento estándar PyTorch en un rango de
6 learning rates logarítmicamente espaciados ($6.25 \times 10^{-5}$, $2.5 \times 10^{-4}$,
$1 \times 10^{-3}$, $4 \times 10^{-3}$, $1.6 \times 10^{-2}$, $6.4 \times 10^{-2}$). Ambos
métodos entrenan durante 10 epochs.

\begin{table}[htbp]
  \centering
  \caption{Precisión de \texttt{qsimov\_gradient} y re-entrenamiento estándar en función del
           learning rate (GPU, PyTorch).}
  \label{tab:tabla3}
  \begin{tabular}{ccc}
    \toprule
    \textbf{Learning rate} & \textbf{\texttt{qsimov\_gradient} (\%)} & \textbf{\texttt{standard} PyTorch (\%)} \\
    \midrule
    $6.25 \times 10^{-5}$ & 94.3 & 87.3 \\
    $2.5  \times 10^{-4}$ & 97.5 & 96.4 \\
    $1    \times 10^{-3}$ & 97.9 & 97.7 \\
    $4    \times 10^{-3}$ & 98.1 & 98.2 \\
    $1.6  \times 10^{-2}$ & 98.2 & 97.0 \\
    $6.4  \times 10^{-2}$ & 97.8 & 44.4 \\
    \bottomrule
  \end{tabular}
\end{table}

La red estándar degrada significativamente con el learning rate más alto (44.4\,\%), mientras que
\texttt{qsimov\_gradient} se mantiene en el rango 94--98\,\% en todo el espectro. Esta robustez se
explica porque la capa \texttt{CustomConnectedLinear} del módulo de gradiente tiene muchas menos
conexiones que la capa Dense original (solo los caminos activos), lo que reduce el espacio de
búsqueda y estabiliza el entrenamiento ante tasas de aprendizaje elevadas.

\subsection{Experimento 3: Olvido catastrófico (MNIST)}

La velocidad y la robustez al learning rate quedan establecidas en los experimentos anteriores. La
pregunta central del TFG es si Qsimov resuelve el problema del olvido catastrófico. Este
experimento lo verifica directamente: un modelo entrenado sobre los dígitos 0--4 se actualiza con
datos exclusivamente de los dígitos 5--9, sin acceso a los datos originales.

\paragraph{Configuración} MNIST dividido en dos fases: fase 1 con los dígitos 0--4 (30.596 train,
5.139 test), fase 2 con los dígitos 5--9 (29.404 train, 4.861 test). Ejecutado con PyTorch sobre
GPU. El modelo base se pre-entrena sobre la fase 1 durante 10 epochs. En la fase 2 no se permite
acceder a los datos de la fase 1 (excepto en el método \texttt{standard\_cumulative}, que sirve de
cota superior teórica).

\begin{table}[htbp]
  \centering
  \caption{Métodos comparados en el experimento de olvido catastrófico (MNIST).}
  \label{tab:metodos_forgetting_mnist}
  \begin{tabular}{lcp{7cm}}
    \toprule
    \textbf{Método} & \textbf{Acc. fase 1} & \textbf{Descripción} \\
    \midrule
    \texttt{base}              & ---        & Modelo original, sin actualización (referencia) \\
    \texttt{linear\_accum}     & No         & QsimovLinearSystem: fit fase 1 + fit fase 2 sin reset \\
    \texttt{linear\_new\_only} & No         & QsimovLinearSystem: solo datos de fase 2 (reset entre fases) \\
    \texttt{gradient\_new\_only} & No       & QsimovGradient: 5 epochs solo en datos de fase 2 \\
    \texttt{finetune\_new\_only} & No       & Fine-tuning estándar PyTorch: 5 epochs solo en datos de fase 2 \\
    \texttt{cumulative}        & \textbf{Sí} & Fine-tuning estándar con datos de fase 1 + fase 2 (acumulado) \\
    \bottomrule
  \end{tabular}
\end{table}

\begin{table}[htbp]
  \centering
  \caption{Resultados del experimento de olvido catastrófico en MNIST (GPU, PyTorch).}
  \label{tab:tabla4}
  \small
  \begin{tabular}{lcccc}
    \toprule
    \textbf{Método} & \textbf{Acc. 0--4 (\%)} & \textbf{Acc. 5--9 (\%)} & \textbf{Acc. global (\%)} & \textbf{Tiempo (s)} \\
    \midrule
    \texttt{base}              &  99.6        &  0.0         & 51.2        & 0.0 \\
    \texttt{linear\_accum}     & \textbf{94.3} & \textbf{49.1} & \textbf{72.3} & \textbf{1.2} \\
    \texttt{linear\_new\_only} &   0.0        & 68.5         & 33.3        & 1.1 \\
    \texttt{gradient\_new\_only} & 0.0        & 61.1         & 29.7        & 7.7 \\
    \texttt{finetune\_new\_only} & 0.0        & 96.1         & 46.7        & 6.4 \\
    \texttt{cumulative}        &  98.5        & 94.8         & 96.7        & 12.6 \\
    \bottomrule
  \end{tabular}
\end{table}

\texttt{linear\_accum} es el único método que mantiene precisión en las clases antiguas (94.3\,\%)
mientras aprende las nuevas (49.1\,\%), y lo hace en \textbf{1.2\,s}, diez veces más rápido que
el re-entrenamiento acumulado y sin acceder a los datos de la fase 1. Los métodos
\texttt{linear\_new\_only}, \texttt{gradient\_new\_only} y \texttt{finetune\_new\_only} exhiben
olvido catastrófico total en las clases 0--4 (0\,\%). El re-entrenamiento acumulado
(\texttt{cumulative}) consigue el mejor rendimiento global (96.7\,\%) pero requiere acceso a los
datos históricos y tarda 12.6\,s. La precisión del \texttt{linear\_accum} en clases nuevas (49.1\,\%)
es inferior a la del fine-tuning sobre nuevas (96.1\,\%): este compromiso es el precio de no olvidar,
ya que el sistema lineal debe satisfacer simultáneamente las ecuaciones de las 30.596 muestras de la
fase 1 y las 29.404 de la fase 2.

\subsection{Experimento 4: CIFAR-10. Olvido catastrófico y datos escasos}

\paragraph{Configuración} Dataset CIFAR-10 (50.000 muestras de entrenamiento, 10.000 de test;
imágenes en color $32 \times 32$ píxeles~\cite{ref18}). Red base: LeNet adaptado con dos bloques
Conv2D+MaxPooling y tres capas Dense. Ejecutado con PyTorch sobre GPU. Punto de corte:
\texttt{initial\_layer = -1}. Número de caminos generados: \textbf{33}.

\textbf{4a.~Olvido catastrófico}: Modelo pre-entrenado en clases 0--4 (fase 1). Actualización con
clases 5--9 (fase 2) sin acceso a datos históricos. 20 epochs en fase 1, 5 epochs en fase 2.

\begin{table}[htbp]
  \centering
  \caption{Resultados del experimento de olvido catastrófico en CIFAR-10 (GPU, PyTorch).}
  \label{tab:tabla5}
  \resizebox{\textwidth}{!}{%
  \begin{tabular}{lcccc}
    \toprule
    \textbf{Método} & \textbf{Acc. antiguas (\%)} & \textbf{Acc. nuevas (\%)} & \textbf{Acc. global (\%)} & \textbf{Tiempo (s)} \\
    \midrule
    \texttt{base} (sin actualizar)   & 78.6         & 0.0           & 39.3        & 0.0 \\
    \texttt{linear\_accum}           & \textbf{73.4} & 8.7          & 41.1        & \textbf{1.3} \\
    \texttt{linear\_new\_only}       & 0.0          & 52.8          & 26.4        & 1.3 \\
    \texttt{gradient\_new\_only}     & 0.0          & 52.4          & 26.2        & 7.7 \\
    \texttt{finetune\_new\_only}     & 0.0          & \textbf{85.2} & 42.6        & 5.7 \\
    \texttt{cumulative} (acumulado)  & 62.7         & 72.9          & \textbf{67.8} & 11.6 \\
    \bottomrule
  \end{tabular}}%
\end{table}

\texttt{finetune\_new\_only} exhibe olvido catastrófico total: la precisión en las clases antiguas
cae de 78.6\,\% a 0.0\,\%. \texttt{linear\_accum} es el único método que preserva conocimiento de
las clases antiguas (73.4\,\%) sin acceso a datos históricos, y lo hace en solo 1.3\,s. La
precisión en clases nuevas es baja (8.7\,\%) porque con solo 33 caminos el espacio representacional
del sistema lineal es muy reducido: resolver simultáneamente las ecuaciones de la fase 1 y de la
fase 2 deja pocos grados de libertad para las clases nuevas. Los métodos
\texttt{linear\_new\_only} y \texttt{gradient\_new\_only} obtienen $\sim$52\,\% en clases nuevas
pero olvidan completamente las antiguas, comportándose como el fine-tuning estándar.

\textbf{4b.~Generalización con datos escasos}: Mismo modelo LeNet entrenado con distintos
subconjuntos del dataset completo (100, 1.000 y 5.000 muestras), comparando gradiente estándar
versus \texttt{qsimov\_gradient}. 10 epochs en todos los casos.

\begin{table}[htbp]
  \centering
  \caption{Precisión en test según el tamaño del split de entrenamiento (GPU, PyTorch, 10 epochs).}
  \label{tab:tabla5b}
  \small
  \begin{tabular}{p{4cm}cc}
    \toprule
    \textbf{Tamaño del split} & \textbf{Std. gradient (\%)} & \textbf{Qsimov gradient (\%)} \\
    \midrule
    100 muestras              & 10.7          & \textbf{20.0} \\
    1.000 muestras            & 25.7          & \textbf{29.8} \\
    5.000 muestras            & 36.8          & 36.0 \\
    Dataset completo (50.000) & \textbf{54.9} & --- \\
    \bottomrule
  \end{tabular}
\end{table}

Con conjuntos de datos muy pequeños (100 muestras), \texttt{qsimov\_gradient} supera al gradiente
estándar en $\sim$9 puntos porcentuales (20.0\,\% vs.\ 10.7\,\%). Esta ventaja se reduce con más
datos y prácticamente desaparece en torno a las 5.000 muestras (36.8\,\% vs.\ 36.0\,\%), donde la
mayor capacidad del gradiente estándar sobre toda la red compensa la regularización implícita de
Qsimov. El resultado confirma que la restricción estructural de Qsimov es una regularización
implícita efectiva en el régimen de escasez de datos.

\subsection{Experimento 5: Aprendizaje continual en ImageNet (100 clases)}

\paragraph{Configuración} Subconjunto de 100 clases de ImageNet (130.000 imágenes de entrenamiento,
$\sim$13.000 de test). Red base: VGG16~\cite{ref19} con la última capa adaptada a 100 clases
(pre-entrenado sobre estas 100 clases desde cero). El escenario de aprendizaje continual divide los
datos en 4 rondas de 25 clases nuevas cada una. Los cuatro métodos comparados fueron
\texttt{qsimov\_linear\_accum}, \texttt{qsimov\_gradient}, \texttt{standard\_finetune} y
\texttt{standard\_cumulative}.

\begin{table}[htbp]
  \centering
  \caption{Resultados globales del experimento de aprendizaje continual sobre ImageNet (100 clases).}
  \label{tab:tabla6}
  \begin{tabular}{lcc}
    \toprule
    \textbf{Método} & \textbf{Acc. global (\%)} & \textbf{Tiempo total (s)} \\
    \midrule
    \texttt{qsimov\_linear\_accum}       & \textbf{37.0} & \textbf{113} \\
    \texttt{standard\_cumulative} (acum.)& 23.0          & 655 \\
    \texttt{standard\_finetune}          & 20.8          & --- \\
    \texttt{qsimov\_gradient}            &  2.5          & --- \\
    \bottomrule
  \end{tabular}
\end{table}

Qsimov con acumulación de ecuaciones supera al re-entrenamiento acumulado en precisión global
(37.0\,\% vs.\ 23.0\,\%) y lo hace seis veces más rápido (113\,s vs.\ 655\,s). La preservación del
conocimiento entre rondas es notable: la precisión en las rondas anteriores se mantiene entre el
79\,\% y el 84\,\% a lo largo de todo el experimento.

El módulo de gradiente de Qsimov no converge en este escenario (2.5\,\% global): con 100 clases y
un learning rate reducido, el método iterativo requeriría muchas más epochs de las utilizadas. Esta
es una limitación identificada del componente de gradiente en problemas de alta complejidad.

\section{Resultados y discusión}

Los cinco experimentos permiten extraer las siguientes conclusiones cuantitativas:

\begin{enumerate}
  \item El sistema lineal de Qsimov elimina el olvido catastrófico en todos los escenarios
        evaluados. En MNIST, \texttt{linear\_accum} mantiene 94.3\,\% en clases antiguas mientras
        aprende un 49.1\,\% de las nuevas, en solo 1.2\,s. En CIFAR-10, preserva 73.4\,\% en
        clases antiguas frente al 0.0\,\% del fine-tuning estándar, en 1.3\,s. En ImageNet (100
        clases, 4 rondas), termina con 37.0\,\% de precisión global, superando al re-entrenamiento
        acumulado (23.0\,\%) y a una velocidad $6\times$ superior.

  \item La elección del \texttt{initial\_layer} condiciona la capacidad del sistema. En
        MNIST, \texttt{initial\_layer = -2} restringe el re-entrenamiento a la última capa
        \texttt{Linear(16} $\rightarrow$ \texttt{10)}, generando 17 caminos; pese a ser pocos, resultan
        suficientes para una tarea relativamente simple (dígitos $28 \times 28$) y el sistema
        alcanza un 49.1\,\% en clases nuevas. En CIFAR-10, \texttt{initial\_layer = -1} genera 33
        caminos para un problema más complejo (imágenes naturales en color), pero ese espacio
        representacional resulta insuficiente: resolver simultáneamente las ecuaciones de la fase 1
        y las de la fase 2 deja pocos grados de libertad para las clases nuevas (8.7\,\%). La
        idoneidad del punto de corte depende, por tanto, de la relación entre el número de caminos
        generados y la complejidad intrínseca del problema, lo que convierte al
        \texttt{initial\_layer} en el hiperparámetro más crítico del sistema.

  \item El sistema lineal es más rápido que el re-entrenamiento estándar tanto en escenarios de re-entrenamiento puro ($\times 22$ en MNIST respecto al
        entrenamiento estándar completo; $\times 9$ en CIFAR-10 respecto al re-entrenamiento
        acumulado) como en aprendizaje continual ($\times 10$ en MNIST respecto al re-entrenamiento
        acumulado; $\times 6$ en ImageNet).

  \item El módulo de gradiente es una regularización implícita efectiva en regímenes de
        escasez de datos: con 100 muestras en CIFAR-10, \texttt{qsimov\_gradient} supera al
        gradiente estándar en $\sim$9 puntos (20.0\,\% vs.\ 10.7\,\%). La ventaja desaparece con
        conjuntos grandes donde la mayor capacidad del gradiente estándar domina (5.000 muestras:
        36.8\,\% vs.\ 36.0\,\%).

  \item El módulo de gradiente es robusto al learning rate en problemas de baja
        complejidad (MNIST: 94--98\,\% en todo el rango de LRs evaluado) pero no escala bien a
        problemas de alta complejidad (ImageNet, 100 clases: 2.5\,\% global con
        LR$=1\times 10^{-5}$ y 10 epochs).
\end{enumerate}

\section{Evaluación presupuestaria}

\begin{table}[htbp]
  \centering
  \caption{Evaluación presupuestaria: estimado vs.\ real.}
  \label{tab:presupuesto}
  \begin{tabular}{lcc}
    \toprule
    \textbf{Concepto} & \textbf{Estimado} & \textbf{Real} \\
    \midrule
    Horas de desarrollo           & 475\,h    & $\sim$490\,h \\
    Coste de personal (15\,€/h)   & 7.125\,€  & 7.350\,€ \\
    Coste GPU (cloud)             & 40\,€     & $\sim$55\,€ \\
    \midrule
    \textbf{Total}                & \textbf{7.165\,€} & \textbf{$\sim$7.405\,€} \\
    \bottomrule
  \end{tabular}
\end{table}

La desviación respecto a la estimación es de un 3,4\,\%, dentro del margen habitual para proyectos
de esta naturaleza.
